\section{Phase 2: unmixing three pairs}\label{sec:fastp2}

Phase~1 hands us two things: the twin
pairings, and the pure columns. However, this is insufficient for Sidelnikov--Shestakov~\cite{sidelnikov1992insecurity}, 
as that requires $k+3$ pure columns, and at the $w=n-k$ parameters, Phase~1 finds exactly $k$
pure columns. 

A trivial method of finding the pairs would be to compute linear combination $c_{1,i}+a_i*c_{2,i}$ for $a \neq 0$
 for each of 3 
pairs $(c_{1,i}, c_{2,i})$ for which we wish to recover pure columns. Exactly 1 of the $(q-1)^3$ possibilities of 
choices of 3 $a_i$ will 
successfully unmix all the 3 columns, enabling the Sidelnikov--Shestakov attack to succeed.

Instead, we use some clever ideas to \emph{solve} for the triple.
It costs only about $q^2$ cheap steps.

\subsection{Putting the pairs in systematic form}
We begin with the matrix \[G = [ P \mid c_{1,1} \mid c_{1,2} \mid c_{2,1} \mid c_{2,2} \mid c_{3,1} \mid c_{3,2}]\] where
 $P$ is the $k \times k$
matrix of the pure columns. After
a left multiply of $P^{-1}$, the generator looks like
$G=[I_k \mid \mathbf{g}_1\ \mathbf{g}_2\ \mathbf{g}_3]$,
where for pair $i$ with public columns $c_{i,1}, c_{i,2}$ we set
$\mathbf{a}_i = P^{-1} c_{i,1}$,
$\mathbf{b}_i = P^{-1} c_{i,2}$, and
\[
    \mathbf{g}_i(\lambda_i) \;=\; \mathbf{a}_i + \lambda_i \mathbf{b}_i,
    \qquad \lambda_i \in \mathbb{P}^1(\Fq),
\]
with $\mathbf{g}_i(\infty) := \mathbf{b}_i$. 

\subsection{The Test}

The dual $D(\lambda)$ is generated by the rows of
\[
    H \;=\; \bigl[\, -\mathbf{g}_1^{T}\ ;\ -\mathbf{g}_2^{T}\ ;\
    -\mathbf{g}_3^{T} \;\big|\; I_3 \,\bigr]
\]
 For $j \in [k]$
write $\gamma_j = (\mathbf{g}_1[j], \mathbf{g}_2[j],
\mathbf{g}_3[j]) \in \Fq^3$ for the $j$-th column of the left
block, and
\[
    \mathbf{m}_j \;=\; \bigl(\gamma_{j1}\gamma_{j2},\
    \gamma_{j1}\gamma_{j3},\ \gamma_{j2}\gamma_{j3}\bigr) \in \Fq^3,
    \qquad
    M^{*}(\lambda) := \bigl[\mathbf{m}_1 \mid \cdots \mid
    \mathbf{m}_k\bigr] \in \Fq^{3\times k}.
\]
In words, $\mathbf{m}_j$ is just the three pairwise products of
the entries of $\gamma_j$. The test is about the rank of this
$3 \times k$ matrix.

\begin{lemma}[Dual-square characterization]\label{lem:fastp2-dual}
Assume $k \geq 2$. Then
\[
    \dim D(\lambda)^{2} \;=\; 3 + \rk M^{*}(\lambda),
\]
and $C(\lambda)$ is a GRS code with distinct support if and only if
\begin{enumerate}
\item[(i)] $\rk M^{*}(\lambda) = 2$;
\item[(ii)] the (then unique up to scale) vector $\mathbf{c} \in
\Fq^3$ with $\mathbf{c}\, M^{*}(\lambda) = 0$ satisfies $c_{12}
c_{13} c_{23} \neq 0$;
\item[(iii)] the $k+3$ columns of $H$ are pairwise
non-proportional.
\end{enumerate}
\end{lemma}

\begin{proof}
Pairwise products of the three generators of $D(\lambda)$ span
$D(\lambda)^2$. The evaluation map sending a quadric
$Q \in \mathrm{Sym}^2(\Fq^3)$, a space of dimension $6$, to the
vector of its values on the columns of $H$ has as its kernel
exactly the quadrics vanishing at all $k+3$ columns of $H$,
viewed as points of $\mathbb{P}^2$. Vanishing at
$\mathbf{e}_1, \mathbf{e}_2, \mathbf{e}_3$ forces the squared
coefficients of $Q$ to be zero, so every quadric in the kernel
looks like $Q_{\mathbf{c}} = c_{12}X_1X_2 + c_{13}X_1X_3 +
c_{23}X_2X_3$, and $Q_{\mathbf{c}}(\gamma_j) =
\mathbf{c}\cdot\mathbf{m}_j$. Hence the kernel has dimension
$3 - \rk M^{*}$, and $\dim D(\lambda)^2 = 6 - (3 - \rk M^{*})
= 3 + \rk M^{*}$, in any characteristic.

If $C(\lambda)$ is GRS with distinct support, so is its dual
$D(\lambda)$~\cite{MS77}, a $\GRS_{k+3,3}$ whose square is the
evaluation code of degree at most $4$ polynomials, of dimension
$\min(5, k+3) = 5$; hence $\rk M^{*} = 2$. The columns of $H$
are the support points of $D(\lambda)$, hence pairwise
non-proportional, and they all lie on a smooth conic. That conic
is a quadric through $\mathbf{e}_1,\mathbf{e}_2,\mathbf{e}_3$,
so it is $Q_{\mathbf{c}}$ for the unique relation
$\mathbf{c}$, and smoothness is exactly (ii): if, say,
$c_{12} = 0$ then $Q_{\mathbf{c}} = X_3(c_{13}X_1 + c_{23}X_2)$
is a union of two lines, whereas if all $c_{ij} \neq 0$ the
gradient system in characteristic $2$ has the unique candidate
singular point, the nucleus $\nu = (c_{23}, c_{13}, c_{12})$,
and $Q_{\mathbf{c}}(\nu) = c_{12}c_{13}c_{23} \neq 0$, so $\nu$
is not on the conic and the conic is smooth.  While we do not need 
it for the purposes of this paper, the same conclusion
in odd characteristic follows from the determinant
$c_{12}c_{13}c_{23}/4 \neq 0$.

Conversely, given (i)--(iii), all $k+3$ columns of $H$ lie on
$Q_{\mathbf{c}}$, which is a smooth conic by (ii), and a
$3$-dimensional code whose generator columns are pairwise
distinct points of a smooth conic is a $\GRS_{3}$ code, by
parametrizing the conic by $\mathbb{P}^1$. So $D(\lambda)$ is
GRS and $C(\lambda) = D(\lambda)^{\perp}$ is GRS with distinct
support by (iii).
\end{proof}

\subsection{Solving for the triple, instead of trying it}

Now, we present our root-finding idea. Fix a triple $J = \{j_1,j_2,j_3\} \subset [k]$ of coordinate positions (for simplicity we may take $(0,1,2)$), and let
\[
    \Delta_J(\lambda) \;=\; \det\bigl[\, R_{ab}(j_t) \,\bigr]_{
    (a,b)\,\in\,\{(1,2),(1,3),(2,3)\},\ t\in\{1,2,3\}}
\]
be the corresponding $3\times 3$ minor of $M^{*}(\lambda)$,
where $R_{ab}(j) = \mathbf{g}_a[j]\,\mathbf{g}_b[j]$. Each
$R_{ab}$ is a product of two affine forms, so every term of the
determinant has degree at most $2$ in each $\lambda_i$:
$\Delta_J$ has multidegree at most $(2,2,2)$. By the lemma,
$\Delta_J$ vanishes at the true triple $\lambda^{*}$, and
indeed every $3\times 3$ minor of $M^{*}(\lambda^{*})$ does,
since $M^{*}(\lambda^{*})$ has rank $2$.

Fix $(\lambda_1,\lambda_2)$ and call that a \emph{cell}. The
row $(R_{12}(j))_{j\in J}$ does not involve $\lambda_3$; call
it $\mathbf{e}$. The other two rows are affine in $\lambda_3$:
$(R_{13}(j))_{j\in J} = \mathbf{p} + \lambda_3 \mathbf{q}$ and
$(R_{23}(j))_{j\in J} = \mathbf{r} + \lambda_3 \mathbf{s}$.
Multilinearity of the determinant in the last two rows gives
\[
    \Delta_J \;=\; A\,\lambda_3^{2} + B\,\lambda_3 + C,
\]
with $A, B, C$ each a $3\times 3$ determinant in
$\mathbf{e},\mathbf{p},\mathbf{q},\mathbf{r},\mathbf{s}$.
So at each cell, the admissible $\lambda_3$ are just the roots
of one explicit quadratic, costing about $150$ field operations
per cell, independent of $k$. Solving that quadratic over
$\mathbb{F}_{2^m}$ is closed-form: if $A = 0$ it is linear (or
constant); if $A \neq 0$ and $B = 0$ the unique root is
$\lambda = (C/A)^{2^{m-1}}$; if $A,B \neq 0$, substituting
$\lambda = (B/A)\mu$ gives $\mu^2 + \mu = AC/B^2$, solvable
iff $\operatorname{Tr}(AC/B^2) = 0$, with roots $\mu, \mu+1$.
Heuristically, we have observed that the number of roots over all $q^2$ cells is about $q^2$ in
expectation (with $68{,}406$ observed in a test at $q=256$).

The rest is a cascade, and each stage is cheaper than the last
in the only sense that matters, the number of survivors.
\textbf{Stage~1} computes the roots at every cell.
\textbf{Stage~2} evaluates a second minor $\Delta_{J'}$, for a
disjoint triple $J'$, at each root (about $60$ operations),
leaving about $q$ survivors. \textbf{Stage~3} applies the exact
test of Lemma~\ref{lem:fastp2-dual} to each survivor ($O(k)$
operations each).~\textbf{Stage~4} rebuilds each surviving
candidate in original coordinates and runs SS, accepting only
a triple that also passes the attack's deterministic full-key
verification. The remaining $w - 3$ pairs are then unmixed
against the recovered curve, and the key completed, exactly as
in the original Phase~2 (including the CLT
post-processing~\cite{CLT18}); none of that is changed.

\begin{algorithm}[H]
\caption{Root-finding unmixing (one orientation)}
\label{alg:fastp2}
\KwIn{The $k$ pure columns; three pairs $(c_{i,1},c_{i,2})$; disjoint
triples $J, J' \subset [k]$.}
\KwOut{A verified unmixing triple, or failure.}
$P \leftarrow$ pure-column matrix; compute $P^{-1}$ and
$\mathbf{a}_i, \mathbf{b}_i$ for $i=1,2,3$\;
\For{$(\lambda_1,\lambda_2) \in \Fq^2$}{
    Form $\mathbf{e},\mathbf{p},\mathbf{q},\mathbf{r},\mathbf{s}$ on
    $J$; compute $A,B,C$\;
    \eIf{$A = B = C = 0$}{$\mathcal{R} \leftarrow \{\lambda_3 \in
        \Fq : \Delta_{J'}(\lambda_1,\lambda_2,\lambda_3) =
        0\}$ \tcp*{degenerate cell: scan}}
    {$\mathcal{R} \leftarrow \{\lambda_3 : \Delta_{J'}(\lambda_1,
        \lambda_2, \lambda_3) = 0,\ \lambda_3 \text{ a root of }
        A\lambda_3^2{+}B\lambda_3{+}C\}$\;}
    \For{$\lambda_3 \in \mathcal{R}$}{
        \If{$(\lambda_1,\lambda_2,\lambda_3)$ passes
            Lemma~\ref{lem:fastp2-dual} and SS $+$ verification on
            original coordinates}{\Return
            $(\lambda_1,\lambda_2,\lambda_3)$\;}
    }
}
\Return failure
\end{algorithm}
